\section{X-Işını İkilileri (X-ray Binaries)}

X-ışını ikilileri, bir yıldız ile bir \emph{kompakt cisim}den oluşan ve sistemden güçlü X-ışını ışıması gözlenen ikili yıldız sistemleridir. X-ışını üretiminin temel fiziksel kaynağı, normal yıldızdan kompakt cisme doğru gerçekleşen \emph{kütle aktarımı} ve bu maddenin güçlü kütleçekim alanında yüksek enerjilere ulaşmasıdır.

\subsection{Kompakt Cisim Kavramı}

Kompakt cisim, yıldız evriminin son aşamasında ortaya çıkan, yarıçapına göre son derece büyük kütleye sahip gök cisimleridir. Başlıca kompakt cisim türleri şunlardır:
\begin{itemize}
    \item Beyaz cüceler
    \item Nötron yıldızları
    \item Kara delikler
\end{itemize}

Bir cismin kompaktlığı genellikle Schwarzschild yarıçapı ile ifade edilir:
\begin{equation}
r_s = \frac{2GM}{c^2}
\end{equation}
burada $G$ evrensel kütleçekim sabiti, $M$ cismin kütlesi ve $c$ ışık hızıdır.

Nötron yıldızları için tipik yarıçap $R \sim 10\ \mathrm{km}$ iken kütle $M \sim 1.4\,M_\odot$ mertebesindedir. Kara delikler için olay ufku doğrudan Schwarzschild yarıçapı ile tanımlanır.

\subsection{X-Işını Üretim Mekanizması}

X-ışını ikililerinde ışınımın temel kaynağı, kompakt cisme düşen maddenin kütleçekim potansiyel enerjisinin kinetik enerjiye ve ardından elektromanyetik ışıma enerjisine dönüşmesidir.

Birim kütle başına açığa çıkan potansiyel enerji:
\begin{equation}
\Delta E \approx \frac{GM}{R}
\end{equation}

Bu enerji, tipik yıldız içi nükleer reaksiyonlardan açığa çıkan enerjiden çok daha büyüktür. Örneğin nötron yıldızları için:
\begin{equation}
\frac{GM}{R} \sim 10^{20}\ \mathrm{J\,kg^{-1}}
\end{equation}

Bu nedenle X-ışını ikilileri, galaksideki en parlak X-ışını kaynakları arasında yer alır.

\subsection{Düşük Kütleli X-Işını İkilileri (LMXB)}

Düşük kütleli X-ışını ikilileri (Low-Mass X-ray Binaries, LMXB), kütlesi genellikle
\begin{equation}
M_\star \lesssim 1\,M_\odot
\end{equation}
olan bir yıldız ile bir nötron yıldızı veya kara delikten oluşur.

Bu sistemlerde kütle aktarımı çoğunlukla \emph{Roche lobu taşması} yoluyla gerçekleşir. Roche potansiyeli, dönen ikili sistemlerde etkin potansiyel olarak tanımlanır:
\begin{equation}
\Phi = -\frac{GM_1}{r_1} - \frac{GM_2}{r_2} - \frac{1}{2}\omega^2 r^2
\end{equation}

Yıldız Roche lobunu doldurduğunda, L1 Lagrange noktası üzerinden madde kompakt cisme akar ve bir \emph{akresyon diski} oluşturur.

Akresyon diskinin yerel sıcaklığı yaklaşık olarak:
\begin{equation}
T(r) \approx \left( \frac{3GM\dot{M}}{8\pi \sigma r^3} \right)^{1/4}
\end{equation}
şeklinde ifade edilir. İç bölgelerde sıcaklık $10^7$--$10^8$ K mertebesine ulaşarak X-ışını üretir.

\subsection{Yüksek Eğilimli İkili Sistemler (X-Işını Kelepçeleri)}

Yüksek eğilimli ikili sistemler, yörünge düzlemi gözlemciye neredeyse kenardan görülen ($i \approx 90^\circ$) X-ışını ikilileridir. Bu sistemler literatürde \emph{X-ray dippers} veya \emph{X-ray clamps} olarak adlandırılır.

Bu sistemlerin ayırt edici özelliği, ışık eğrisinde periyodik veya yarı-periyodik X-ışını azalmalarıdır. Bu düşüşler, akresyon diskindeki kalınlaşmış bölgelerin veya gaz akımlarının görüş hattını geçici olarak engellemesiyle oluşur.

X-ışını akısı:
\begin{equation}
F_X = \frac{L_X}{4\pi d^2} e^{-\tau(E)}
\end{equation}
şeklinde yazılabilir. Burada $\tau(E)$ enerjiye bağlı optik derinliktir:
\begin{equation}
\tau(E) = \int n(s)\,\sigma(E)\,ds
\end{equation}

Yüksek eğilimli sistemlerde $\tau(E)$ zamanla değişir ve bu durum X-ışını spektrumunda sertleşme ve yumuşama geçişlerine neden olur.

\subsection{Oluşum Mekanizmaları}

X-ışını ikililerinin oluşumu, ikili yıldız evriminin ileri aşamalarını içerir:
\begin{enumerate}
    \item Başlangıçta iki normal yıldızdan oluşan bir ikili sistem
    \item Daha büyük kütleli yıldızın süpernova patlaması
    \item Geriye nötron yıldızı veya kara delik kalması
    \item Yörüngenin kararlı kalması durumunda kütle aktarımının başlaması
\end{enumerate}

Kütle aktarımı sırasında açısal momentum korunumu:
\begin{equation}
J = \mu \sqrt{GMa}
\end{equation}
ile ifade edilir. Burada $\mu$ indirgenmiş kütle, $a$ yörünge yarıçapıdır. Kütle kaybı ve aktarımı, yörünge periyodunun zamanla evrimleşmesine neden olur.

\subsection{Astrofiziksel Önemi}

X-ışını ikilileri:
\begin{itemize}
    \item Güçlü kütleçekim alanlarının test edilmesi
    \item Genel göreliliğin gözlemsel sınanması
    \item Nötron yıldızlarının denklemi halinin kısıtlanması
    \item Kara delik adaylarının belirlenmesi
\end{itemize}
açılarından modern astrofiziğin temel laboratuvarlarıdır.

Bu sistemler, yüksek enerji astrofiziği ile yıldız evrimi teorileri arasında köprü kurar.

\section{Yüksek Kütleli X-Işını Sistemleri (High-Mass X-ray Binaries)}

Yüksek kütleli X-ışını sistemleri (High-Mass X-ray Binaries, HMXB), bir \emph{kompakt cisim} (nötron yıldızı veya kara delik) ile kütlesi büyük olan ($M_\star \gtrsim 10\,M_\odot$) genç ve sıcak bir yıldızdan oluşan ikili sistemlerdir. Bu sistemlerde X-ışını yayınımı, çoğunlukla eş yıldızın güçlü yıldız rüzgârı veya sınırlı Roche lobu taşması yoluyla gerçekleşen akresyon sonucu ortaya çıkar.

HMXB'ler, yıldız oluşum bölgeleriyle ilişkili olup kısa ömürlü sistemlerdir.

\subsection{Eş Yıldızın Fiziksel Özellikleri}

HMXB'lerde eş yıldız genellikle:
\begin{itemize}
    \item O veya erken B tipi anakol yıldızı
    \item Süperdev yıldız
    \item Be yıldızı (hızlı dönen, diskli)
\end{itemize}

Bu yıldızlar güçlü yıldız rüzgârları üretir. Rüzgâr kütle kaybı hızı tipik olarak:
\begin{equation}
\dot{M}_{\mathrm{wind}} \sim 10^{-7} \text{--} 10^{-5}\, M_\odot\,\mathrm{yr^{-1}}
\end{equation}

Rüzgâr hızları:
\begin{equation}
v_{\mathrm{wind}} \sim 10^3\ \mathrm{km\,s^{-1}}
\end{equation}

\subsection{X-Işını Üretim Mekanizması}

HMXB'lerde X-ışını üretimi, rüzgâr akresyonu veya Roche lobu taşması ile ilişkilidir. Rüzgâr akresyonu için Bondi–Hoyle–Lyttleton yaklaşımı kullanılır:
\begin{equation}
\dot{M}_{\mathrm{acc}} = \frac{4\pi G^2 M^2 \rho}{(v^2 + c_s^2)^{3/2}}
\end{equation}

Burada $M$ kompakt cismin kütlesi, $\rho$ rüzgâr yoğunluğu, $v$ bağıl hız ve $c_s$ ses hızıdır.

X-ışını parlaklığı yaklaşık olarak:
\begin{equation}
L_X \approx \frac{GM\dot{M}_{\mathrm{acc}}}{R}
\end{equation}

\subsection{Yüksek Kütleli X-Işını Sistemlerinin Türleri}

\subsubsection{Süperdev X-Işını Sistemleri (sgHMXB)}

Bu sistemlerde eş yıldız bir süperdevdir. Akresyon doğrudan yıldız rüzgârı yoluyla gerçekleşir. X-ışını parlaklığı genellikle süreklidir:
\begin{equation}
L_X \sim 10^{35} \text{--} 10^{37}\ \mathrm{erg\,s^{-1}}
\end{equation}

Bu sistemlerde X-ışını pulsarları yaygındır.

\subsubsection{Be/X-Işını Sistemleri (BeXRB)}

Bu sistemlerde eş yıldız hızlı dönen bir Be yıldızıdır ve çevresinde ekvatoryal bir gaz diski bulunur. Kompakt cisim eliptik yörüngede hareket eder ve diskten madde yakalar.

X-ışını patlamaları genellikle periyodiktir:
\begin{itemize}
    \item Tip I patlamalar: Yörünge periyodu ile ilişkili
    \item Tip II patlamalar: Disk kararsızlıkları kaynaklı
\end{itemize}

\subsection{Oluşum Süreci}

HMXB'lerin oluşumu şu evreleri içerir:
\begin{enumerate}
    \item İki büyük kütleli yıldızdan oluşan ikili sistem
    \item Daha ağır yıldızın süpernova patlaması
    \item Nötron yıldızı veya kara delik oluşumu
    \item İkili sistemin bağlanmış kalması
    \item Yıldız rüzgârı veya Roche lobu taşması ile akresyon
\end{enumerate}

Bu sistemlerin toplam ömrü genellikle:
\begin{equation}
\tau \sim 10^6\text{--}10^7\ \mathrm{yr}
\end{equation}

\section{Ultra-Parlak X-Işını Sistemleri (Ultra-Luminous X-ray Sources)}

Ultra-parlak X-ışını sistemleri (ULX), galaksi çekirdeği dışında bulunan ve X-ışını parlaklığı Eddington sınırını aşan noktasal kaynaklardır:
\begin{equation}
L_X \gtrsim 10^{39}\ \mathrm{erg\,s^{-1}}
\end{equation}

Bu parlaklık, tipik yıldız kütleli kara delikler için Eddington parlaklığının üzerindedir.

\subsection{Eddington Parlaklığı}

Eddington parlaklığı, radyasyon basıncı ile kütleçekimin dengede olduğu maksimum parlaklıktır:
\begin{equation}
L_{\mathrm{Edd}} = \frac{4\pi GMm_p c}{\sigma_T}
\end{equation}

Burada $m_p$ proton kütlesi, $\sigma_T$ Thomson saçılma kesitidir.

\subsection{ULX'lerin Fiziksel Yorumu}

ULX sistemleri üç ana modelle açıklanır:

\subsubsection{Ara Kütleli Kara Delikler (IMBH)}

Bu senaryoda kompakt cisim:
\begin{equation}
M \sim 10^2\text{--}10^5\,M_\odot
\end{equation}
kütleli bir kara deliktir ve parlaklık Eddington sınırı içinde kalır.

\subsubsection{Süper-Eddington Akresyon}

Yıldız kütleli kara delik veya nötron yıldızı üzerinde, geometrik olarak kalın disk ve ışınım demetlemesi (beaming) söz konusudur:
\begin{equation}
L_{\mathrm{obs}} \approx b^{-1} L_{\mathrm{Edd}}, \quad b < 1
\end{equation}

\subsubsection{ULX Pulsarları}

Bazı ULX'lerin nötron yıldızı içerdiği keşfedilmiştir. Bu sistemlerde:
\begin{itemize}
    \item Güçlü manyetik alan ($B \sim 10^{12}$--$10^{13}$ G)
    \item Kolimasyonlu akresyon
    \item X-ışını pulsasyonları
\end{itemize}
gözlenir.

\subsection{ULX'lerin Astrofiziksel Önemi}

ULX sistemleri:
\begin{itemize}
    \item Süper-Eddington akresyon fiziğinin anlaşılması
    \item Kara delik kütle dağılımının sınırlarının belirlenmesi
    \item Yoğun yıldız oluşum bölgelerinin izlenmesi
    \item Erken evren kuasarlarının analogları
\end{itemize}
olarak yüksek enerji astrofiziğinde kritik bir rol oynar.
